% =============================================================
% VII. CONCLUSION
% =============================================================
\section{Conclusion}
\label{sec:conclusion}

This paper has constructed the first annual, internationally comparable series of Estonian GDP per capita for 1919--95 and joined it to existing estimates to give a continuous series for 1913 and 1919--2022 in 1990 and 2011 international dollars. It rests on a new 1938 benchmark from a direct price comparison with Sweden, a recalculated production-side series for 1919--40, and a new series for the ESSR for 1937--95 built from physical output, transport and employment data and anchored to Bolotin's 1990 benchmark. Built from different sources, with different weights and from benchmarks half a century apart, the two new series agree within 4 per cent in every year of their 1937--40 overlap, and the Soviet-period series meets a modern series derived from the 2011 ICP round within 1 per cent at the 1991 splice. The series supersedes the three partial estimates that preceded it, lowering late-Soviet income by about a quarter relative to the Maddison Project's NMP-based figures and raising it by a similar margin relative to the only earlier long series, and it is ready for the Maddison Project Database.

Three findings stand out. First, Estonia's interwar growth was respectable but not exceptional: \gr{1919}{1929} per cent a year in the 1920s, as the country recovered from war and the loss of its imperial market, and \gr{1929}{1938} per cent in the 1930s, when devaluation and a state-led industrial drive produced a rapid recovery from the Depression. Second, Soviet rule brought two decades of rapid growth and two of stagnation: \gr{1950}{1968} per cent a year in 1950--68, faster than in Finland or the Western core, as investment directed from Moscow built an oil-shale, textile and engineering economy on an unusually literate labour force, and \gr{1968}{1989} per cent in 1968--89, as the sources of extensive growth ran out. Estonia's income relative to the Western core rose from two-thirds in 1950 to seven-tenths in 1968 and fell back to half by 1989, while within the Union it remained what it had been at annexation, a republic between a fifth and a half richer than the average. Third, the transition cut income per head by a third in four years, and recovery after 1995 carried Estonia back, by 2022, close to the share of Western European income it had briefly reached in 1968.

The comparison with Finland, which gives the paper its title, needs restating. The two countries were close in 1913 and 1919, but by 1938 Estonia stood at three-quarters of the Finnish level, and almost half of the gap of 1989 had opened under independence, not occupation; the further divergence was concentrated after 1968, Estonia having kept pace with Finland in the preceding two decades. Finland shows what institutions can do with a given endowment of human capital, but not what Estonia would have become, because the Finnish path depended on resources and markets that Estonia did not have. The costs of Soviet rule were real, and the series locates most of them in the stagnation of 1968--89; measured against the Western core or periphery rather than against a Finland whose growth in those decades was exceptional, they are smaller than the Finnish comparison implies, and they are not the costs of a Finnish miracle that Estonia was denied.

Several extensions are open: a dedicated study of the 1938 service sector of the kind now available for Lithuania and Latvia; a Soviet-period benchmark built from the Estonian archives with republican rather than Union-wide prices, which would test the level to which the volume indices are anchored; expenditure-side cross-checks for a few interwar years; and the application of the procedure to the other non-Russian republics, whose national accounts remain the largest gap in the quantitative economic history of the Soviet Union.
